% Manual de Lázaro Linux
% Compilar:  ./compilar.sh   (o: latexmk lazaro-manual.tex)
% Licencia del texto y las imágenes: CC BY-SA 4.0 (ver LICENCIA.md)
\documentclass[12pt,a4paper,openany]{book}
% Estilo del Manual de Lázaro Linux.
% Aquí solo va la apariencia; el texto está en capitulos/.

% --- Idioma y letras -------------------------------------------------------
\usepackage[T1]{fontenc}
\usepackage[utf8]{inputenc}
\usepackage[spanish,es-noshorthands,es-tabla]{babel}
\usepackage{tgpagella}            % texto: Pagella, fácil de leer en pantalla y papel
\usepackage{tgheros}              % títulos y recuadros: Heros (sin remates)
\renewcommand{\ttdefault}{lmtt}  % comandos de la terminal: Latin Modern Mono
\usepackage{microtype}
\usepackage{amssymb}
\linespread{1.08}

% --- Página ---------------------------------------------------------------
\usepackage[a4paper,top=2.4cm,bottom=2.6cm,inner=2.6cm,outer=2.4cm,headheight=15pt]{geometry}
\usepackage[dvipsnames]{xcolor}
\usepackage{graphicx}
\graphicspath{{imagenes/}{imagenes/capturas/}{../live/usr/share/calamares/branding/lazaro/}}
\usepackage{booktabs,tabularx,array}
\usepackage{enumitem}
\setlist{itemsep=0.25em}
\usepackage{fancyhdr}
\usepackage{titlesec}
\usepackage{tcolorbox}
\tcbuselibrary{skins,breakable}
\usepackage{csquotes}
\usepackage{menukeys}             % \menu{Archivo > Guardar}, \keys{\ctrl + C}
\usepackage[backend=biber,style=numeric,sorting=none,refsegment=chapter,defernumbers=true]{biblatex}
\addbibresource{referencias.bib}
\usepackage{hyperref}

% --- Colores de Lázaro ----------------------------------------------------
\definecolor{pulso}{HTML}{C62828}     % rojo del logo (el pulso)
\definecolor{tinta}{HTML}{263238}
\definecolor{consejo}{HTML}{2E7D32}
\definecolor{cuidado}{HTML}{E65100}
\definecolor{nota}{HTML}{1565C0}

\hypersetup{
  colorlinks=true, linkcolor=tinta, urlcolor=nota, citecolor=consejo,
  pdftitle={Manual de Lázaro Linux},
  pdfauthor={Proyecto Lázaro Linux},
  pdfsubject={Guía de uso de Lázaro Linux para todos},
  pdflang={es},
  bookmarksnumbered=true,
}

% --- Títulos ----------------------------------------------------------------
\titleformat{\chapter}[display]
  {\sffamily\bfseries\color{tinta}}
  {\color{pulso}\Large Capítulo \thechapter}
  {0.4em}
  {\Huge}
  [\vspace{0.3em}{\color{pulso}\titlerule[1.5pt]}]
\titlespacing*{\chapter}{0pt}{-1.5em}{1.8em}
\titleformat{\section}{\sffamily\bfseries\Large\color{tinta}}{\color{pulso}\thesection}{0.8em}{}
\titleformat{\subsection}{\sffamily\bfseries\large\color{tinta}}{\thesubsection}{0.8em}{}
\titleformat{\part}[display]{\centering\sffamily\bfseries\color{tinta}}
  {\color{pulso}\huge Parte \thepart}{0.6em}{\Huge}

\pagestyle{fancy}
\fancyhf{}
\renewcommand{\headrulewidth}{0.4pt}
\fancyhead[LE]{\sffamily\small\leftmark}
\fancyhead[RO]{\sffamily\small\rightmark}
\fancyfoot[C]{\sffamily\small\thepage}
\fancypagestyle{plain}{\fancyhf{}\renewcommand{\headrulewidth}{0pt}\fancyfoot[C]{\sffamily\small\thepage}}

% --- Recuadros --------------------------------------------------------------
% \begin{consejo} ... \end{consejo}   algo que ayuda
% \begin{cuidado} ... \end{cuidado}   algo que puede salir mal
% \begin{nota} ... \end{nota}         un dato extra
\newcommand{\recuadro}[3]{%
  \newtcolorbox{#1}{enhanced,breakable,colback=#2!6,colframe=#2,
    boxrule=0pt,leftrule=3pt,arc=0pt,outer arc=0pt,
    left=8pt,right=8pt,top=6pt,bottom=6pt,
    fonttitle=\sffamily\bfseries\color{#2},coltitle=#2,colbacktitle=#2!6,
    title={#3},attach title to upper={\par\smallskip}}}
\recuadro{consejo}{consejo}{Consejo}
\recuadro{cuidado}{cuidado}{Cuidado}
\recuadro{nota}{nota}{Para saber más}

% Pasos numerados grandes: \begin{pasos} \item ... \end{pasos}
\newcommand{\numpaso}[1]{\tcbox[on line,colback=pulso,colframe=pulso,coltext=white,
  boxsep=1pt,left=3pt,right=3pt,top=1pt,bottom=1pt,arc=2pt]{\sffamily\bfseries #1}}
\newlist{pasos}{enumerate}{1}
\setlist[pasos]{label=\protect\numpaso{\arabic*},leftmargin=2.2em,itemsep=0.5em}

% Lo que se escribe en la terminal.
\newtcolorbox{terminal}{enhanced,colback=tinta,colupper=white,boxrule=0pt,arc=2pt,
  left=8pt,right=8pt,top=5pt,bottom=5pt,fontupper=\ttfamily\footnotesize}
\newcommand{\comando}[1]{\texttt{#1}}

% --- Capturas de pantalla ----------------------------------------------------
% \captura{archivo}{pie de foto}
% Si la imagen todavía no existe, deja un recuadro gris para no frenar la
% escritura: el PDF se compila igual y se ve qué capturas faltan.
\newcommand{\captura}[3][0.85\textwidth]{%
  \begin{figure}[htbp]\centering
  \IfFileExists{imagenes/capturas/#2.png}{%
    \fcolorbox{black!25}{white}{\includegraphics[width=#1]{#2}}}{%
    \IfFileExists{imagenes/capturas/#2.jpg}{%
      \fcolorbox{black!25}{white}{\includegraphics[width=#1]{#2}}}{%
      \fcolorbox{black!30}{black!5}{\parbox[c][0.28\textheight][c]{#1}{\centering
        \sffamily\color{black!55}Captura pendiente:\\\texttt{#2}}}}}
  \caption{#3}\label{fig:#2}
  \end{figure}}

% Términos del glosario: \termino{UEFI} pone la palabra en cursiva la
% primera vez; el glosario completo está en el apéndice.
\newcommand{\termino}[1]{\emph{#1}}

% Referencias al final de cada capítulo: \fuentes imprime solo las citadas
% en ese capítulo. La lista completa va en el apéndice.
\newcommand{\fuentes}{\printbibliography[segment=\therefsegment,
  heading=subbibliography,title={Fuentes y lecturas del capítulo}]}

\input{version}% la crea compilar.sh con la versión de config.env
\setcounter{tocdepth}{1}

\begin{document}
\frontmatter
\begin{titlepage}
\centering
\vspace*{2cm}
\includegraphics[width=5cm]{logo}\par
\vspace{1.5cm}
{\sffamily\bfseries\fontsize{38}{44}\selectfont\color{tinta} Manual de\\[0.2em] Lázaro Linux\par}
\vspace{0.8cm}
{\sffamily\Large\color{pulso} Tu hardware aún tiene pulso\par}
\vspace{1.5cm}
{\large Guía para instalar y usar Lázaro Linux en el día a día,\\
sin conocimientos técnicos\par}
\vfill
{\sffamily Versión \lazaroversion\ \textperiodcentered\ \lazarofecha\par}
\end{titlepage}

\thispagestyle{empty}
\vspace*{\fill}
\noindent{\small
\textbf{Manual de Lázaro Linux}, versión \lazaroversion.\\
Proyecto Lázaro Linux, \url{https://github.com/Entropio1984/lazaro-linux}.

\medskip
\noindent Este manual es libre: puedes copiarlo, imprimirlo, compartirlo y
mejorarlo, siempre que digas de dónde viene y compartas tus cambios con la
misma licencia, la Creative Commons Atribución-CompartirIgual 4.0
Internacional (CC BY-SA 4.0):
\url{https://creativecommons.org/licenses/by-sa/4.0/deed.es}.

\medskip
\noindent Los nombres de otros programas y marcas que aparecen aquí son de sus
dueños y se mencionan solo para explicar cómo usarlos.

\medskip
\noindent ¿Encontraste un error o falta algo? Avísanos en
\url{https://github.com/Entropio1984/lazaro-linux/issues}.\par}
\cleardoublepage

\tableofcontents
\chapter*{Cómo usar este manual}
\addcontentsline{toc}{chapter}{Cómo usar este manual}
\markboth{Cómo usar este manual}{Cómo usar este manual}

Este manual es para cualquier persona que quiera usar Lázaro Linux, aunque
nunca haya usado otra cosa que Windows o el celular. No hace falta saber
de computación: cada palabra técnica se explica la primera vez que
aparece, y todas están juntas en el \hyperref[cap:a-glosario]{Glosario} del
final.

\paragraph{No hace falta leerlo entero.} Está dividido en cuatro partes:
\begin{itemize}
  \item \textbf{Empezar} (capítulos 1 a 5): qué es Lázaro, cómo se prepara el
        USB y cómo se instala. Léela de corrido la primera vez.
  \item \textbf{Uso diario} (capítulos 6 a 14): internet, documentos, PDF,
        videos, impresoras, programas. Ve directo a lo que necesites.
  \item \textbf{Herramientas de Lázaro} (capítulo 15): los programas propios
        de Lázaro para que el equipo vaya mejor.
  \item \textbf{Cuando algo falla} (capítulos 16 a 18): qué hacer si algo
        no funciona y cómo pedir ayuda.
\end{itemize}

\paragraph{Cómo se escriben las cosas.}
\begin{itemize}
  \item Los botones, menús y opciones que ves en la pantalla van así:
        \menu{Archivo > Guardar como}. Quiere decir: abre el menú
        \emph{Archivo} y elige \emph{Guardar como}.
  \item Las teclas van así: \keys{\ctrl + C} quiere decir \enquote{mantén
        apretada la tecla Ctrl y, sin soltarla, pulsa la C}.
  \item Lo que se escribe en la terminal va en un recuadro oscuro. No hace
        falta usar la terminal para casi nada, pero a veces es el camino más
        corto:
\end{itemize}
\begin{terminal}
lazaro-afinador --revisar
\end{terminal}

\paragraph{Los recuadros de colores.}
\begin{consejo}
Un truco que te ahorra tiempo.
\end{consejo}
\begin{cuidado}
Algo que puede hacerte perder archivos o dejar el equipo sin arrancar.
Léelo siempre.
\end{cuidado}
\begin{nota}
Un dato extra para quien quiera entender más. Puedes saltarlo.
\end{nota}

\paragraph{Las fuentes.} Al final de cada capítulo hay una lista de páginas
de confianza (la documentación oficial de Alpine, de cada escritorio y de
cada programa) para profundizar. Los números entre corchetes en el texto,
como \cite{alpine-about}, llevan a esa lista. Todas juntas están en la
Bibliografía, al final del libro.

\paragraph{Escritorios distintos.} Lázaro ofrece varios escritorios (el
\enquote{aspecto} del sistema). Las capturas de pantalla son de
\textbf{XFCE}, el que viene marcado de fábrica. Si elegiste otro, los
nombres de los menús pueden cambiar un poco; el capítulo~\ref{cap:06-escritorio}
explica las diferencias.


\mainmatter
\part{Empezar}
\chapter{Bienvenida}\label{cap:01-bienvenida}

\section{Qué es Lázaro Linux}

Lázaro Linux es un \termino{sistema operativo}: el programa principal de la
computadora, el que se enciende primero y deja usar todos los demás. Cumple
el mismo papel que Windows o macOS. Es gratuito, se instala desde una
memoria USB y está pensado para darle una segunda vida a computadoras y
portátiles viejos que con Windows ya van demasiado lentos o que dejaron de
recibir actualizaciones.

Con Lázaro puedes, sin instalar nada más:
\begin{itemize}
  \item navegar por internet con Firefox;
  \item escribir cartas, hacer planillas y presentaciones con LibreOffice, y
        abrir y guardar archivos de Word, Excel y PowerPoint;
  \item abrir, leer, rellenar e imprimir archivos PDF;
  \item ver YouTube sin que el video se trabe, incluso en equipos muy lentos;
  \item escuchar música, ver películas y fotos;
  \item imprimir y escanear;
  \item hacer videollamadas desde el navegador.
\end{itemize}

El nombre viene del Lázaro de la historia que vuelve a la vida. El lema,
\emph{Tu hardware aún tiene pulso}, quiere decir eso: tu equipo todavía
sirve.

\section{Linux, Alpine y Lázaro}

\termino{Linux} es una familia de sistemas operativos libres. Hay muchos,
llamados \termino{distribuciones}: Ubuntu, Linux Mint, Fedora, Debian y
cientos más. Todos comparten el mismo corazón, el \termino{kernel} Linux,
y se diferencian en los programas que traen y en cómo se ven
\cite{gnu-linux}.

Lázaro está construido sobre \textbf{Alpine Linux}, una distribución
pensada para ser pequeña, simple y segura \cite{alpine-about}. Alpine se
usa sobre todo en servidores, porque gasta muy poca memoria y muy poco
disco. Lázaro toma esa base liviana y le agrega lo que una persona
necesita en su casa: un instalador gráfico en español, un escritorio listo
para usar, los programas de todos los días y unas herramientas propias para
que el equipo funcione bien sin tocar nada.

\begin{nota}
\textbf{Libre} no solo quiere decir gratis. Quiere decir que cualquiera
puede usar el programa para lo que quiera, ver cómo está hecho, mejorarlo
y compartirlo \cite{fsf-software-libre}. Por eso Lázaro puede existir:
reúne programas libres que miles de personas hacen y mejoran en todo el
mundo.
\end{nota}

\section{Qué tiene de distinto}

\begin{description}
  \item[Hecho para equipos viejos.] Al arrancar revisa tu equipo y te
    recomienda el escritorio que mejor le va según la memoria y la tarjeta
    gráfica.
  \item[Sin pasos técnicos.] Idioma, teclado, sonido, wifi, impresoras y
    controladores de video quedan listos solos.
  \item[Funciona sin internet.] La edición Completa trae todo en el USB y
    se instala en un equipo sin red.
  \item[YouTube ligero.] Un reproductor que no usa el navegador y gasta
    mucho menos procesador (capítulo~\ref{cap:10-multimedia}).
  \item[Siempre hay un plan B.] Si el escritorio no llega a abrir, hay un
    instalador en modo texto que hace lo mismo con el teclado.
  \item[Afinado para tu equipo.] El \emph{Afinador} ajusta el sistema para
    que no se congele con muchas pestañas abiertas ni al copiar a un
    pendrive (capítulo~\ref{cap:15-herramientas}).
  \item[Avisa de las actualizaciones.] Un icono en la barra te avisa cuando
    hay novedades, y se instalan con un clic.
\end{description}

\section{¿Funciona en mi equipo?}\label{sec:requisitos}

\begin{center}
\begin{tabularx}{\textwidth}{>{\bfseries}lX}
\toprule
Procesador & De 64 bits (x86\_64). Casi todos los equipos desde 2008 lo son.
             Los muy antiguos de 32 bits (por ejemplo, muchos netbooks con
             Intel Atom N270 o N280) no sirven. \\
Memoria (RAM) & 2 GB o más es lo cómodo. Con 1 GB puede funcionar con el
             escritorio LXQt; con menos, el instalador no deja seguir. \\
Disco & 10 GB como mínimo (el instalador no deja seguir con menos); 32 GB o
             más si vas a guardar fotos y videos. \\
Arranque & BIOS o UEFI: sirven los dos, también los equipos muy viejos. \\
Memoria USB & De 4 GB para la edición Ligera o de 8 GB para la Completa. \\
\bottomrule
\end{tabularx}
\end{center}

Ya se probó en equipos reales como un Lenovo ideapad 110 (Celeron, 4 GB),
un HP Compaq CQ50 (Sempron de un núcleo) y PC con tarjetas NVIDIA de
varias épocas.

\begin{consejo}
¿No sabes cuánta memoria tiene tu equipo? En Windows, pulsa
\keys{Windows + Pausa} o busca \enquote{Acerca de} en el menú Inicio: aparece
como \enquote{RAM instalada}. Ahí mismo dice si el sistema es \enquote{de 64
bits}.
\end{consejo}

\section{Las dos ediciones}\label{sec:ediciones}

Lázaro se descarga en dos versiones, llamadas \emph{ediciones}. Las dos
instalan el mismo sistema; cambia lo que traen dentro del USB.

\begin{center}
\begin{tabularx}{\textwidth}{lXX}
\toprule
 & \textbf{Ligera} & \textbf{Completa} \\
\midrule
Tamaño & menos de 2 GB & unos 4,6 GB \\
¿Necesita internet al instalar? & Para GNOME, KDE Plasma y las aplicaciones
  de Flathub, sí & No (solo para Google Chrome) \\
Escritorios sin internet & XFCE, MATE, LXQt, Sway, XFCE Wayland & Los 7, en
  todos los idiomas \\
Controladores (firmware) & Los de red & Todos \\
\bottomrule
\end{tabularx}
\end{center}

\textbf{Si dudas, elige la Ligera}: se descarga más rápido y alcanza para la
mayoría. Elige la Completa si el equipo no va a tener internet o si tu
conexión es lenta.

\section{Antes de seguir}

Lázaro está todavía en desarrollo. Ya se instala bien en muchos equipos
reales, pero:
\begin{cuidado}
Instalar Lázaro \textbf{borra el disco que elijas}. Antes de instalar,
guarda una copia de tus fotos, documentos y todo lo que te importe en un
pendrive, un disco externo o en la nube. El capítulo~\ref{cap:02-antes-de-instalar}
explica cómo.
\end{cuidado}

Si es tu primera vez, sigue el orden del libro: el capítulo~2 prepara todo,
el~3 te deja probar Lázaro sin tocar nada del equipo, y el~4 lo instala.

\fuentes

\chapter{Antes de instalar}\label{cap:02-antes-de-instalar}

Este capítulo prepara todo lo necesario: una copia de tus archivos, la
descarga de Lázaro y la memoria USB desde la que arrancará el equipo.
Calcula una hora, casi toda esperando a que termine la descarga.

Vas a necesitar:
\begin{itemize}
  \item el equipo donde quieres instalar Lázaro, enchufado a la corriente;
  \item otra computadora con internet para descargar (puede ser la misma,
        si todavía funciona con Windows);
  \item una memoria USB de 4 GB o más (8 GB para la edición Completa) que
        no tenga nada importante: \textbf{se va a borrar}.
\end{itemize}

\section{Guarda una copia de tus archivos}\label{sec:copia-antes}

\begin{cuidado}
Al instalar Lázaro se borra el disco que elijas, con todo lo que tiene. No
hay forma de recuperar lo borrado. Haz la copia \textbf{antes} de seguir.
\end{cuidado}

Lo que suele importar está en estas carpetas de Windows: \emph{Documentos},
\emph{Imágenes}, \emph{Música}, \emph{Videos}, \emph{Descargas} y
\emph{Escritorio}. Cópialas a un pendrive o a un disco externo (no al
mismo pendrive que usarás para Lázaro). Revisa también:
\begin{itemize}
  \item \textbf{Las contraseñas y marcadores del navegador.} Si usas Chrome
        o Edge con una cuenta, ya están guardados en esa cuenta. Si no,
        expórtalos: en Chrome, \menu{Marcadores > Administrador de marcadores
        > \ldots > Exportar marcadores}.
  \item \textbf{El correo.} Si lo lees en una página web (Gmail, Outlook.com,
        Yahoo) no tienes que hacer nada: sigue en internet.
  \item \textbf{Las fotos del celular} que hayas pasado al equipo.
  \item \textbf{Claves de wifi}: anota la de tu casa; la vas a necesitar.
\end{itemize}

Para hacer una copia de seguridad bien hecha y conservarla en el tiempo,
el capítulo~\ref{cap:11-archivos} da más detalles. El Instituto Nacional de
Ciberseguridad de España tiene una buena guía para no expertos
\cite{incibe-copias}.

\section{Descarga Lázaro}\label{sec:descargar}

Primero elige la edición (sección~\ref{sec:ediciones}): \textbf{Ligera} si
tienes buena conexión, \textbf{Completa} si el equipo no va a tener
internet. Luego descarga \emph{dos} archivos de la misma carpeta:

\begin{itemize}
  \item el archivo que termina en \texttt{.iso}: es el \termino{archivo ISO}, un
        \enquote{disco} con todo el sistema dentro. Pesa entre 2 y 4,6 GB;
  \item el archivo del mismo nombre que termina en \texttt{.sha256}: es muy
        pequeño y sirve para comprobar que la descarga llegó bien.
\end{itemize}

Los dos están en:
\begin{itemize}
  \item \textbf{Ligera:} \url{https://github.com/Entropio1984/lazaro-linux/releases}
        o \url{https://sourceforge.net/projects/lazaro-linux/files/Ligera/}
  \item \textbf{Completa:} \url{https://sourceforge.net/projects/lazaro-linux/files/Completa/}
  \item \textbf{Versiones numeradas} (fijas, las más probadas):
        \url{https://sourceforge.net/projects/lazaro-linux/files/Versiones/}
\end{itemize}

No hace falta crear ninguna cuenta. Si la descarga se corta, vuelve a
empezarla: el navegador suele poder continuarla desde donde quedó.

\section{Comprueba que la descarga llegó entera}\label{sec:sha256}

Una descarga dañada es la causa más común de que el USB no arranque, y
tiene arreglo fácil si se detecta a tiempo. El archivo \texttt{.sha256}
trae una \termino{suma de comprobación}: un número largo que se calcula a
partir de cada byte de la ISO. Si cambia un solo byte, cambia el número
entero.

\subsection*{En Windows}
\begin{pasos}
  \item Abre la carpeta donde quedó la descarga (normalmente \emph{Descargas}).
  \item Haz clic derecho en un espacio vacío de la carpeta y elige
        \menu{Abrir en Terminal} (en Windows 10: mantén \keys{\shift}
        apretada mientras haces clic derecho y elige \menu{Abrir la ventana
        de PowerShell aquí}).
  \item Escribe esto y pulsa \keys{Enter}:
\begin{terminal}
Get-FileHash .\textbackslash lazaro-linux-*.iso
\end{terminal}
  \item Tarda un minuto. Aparece una línea con el número (la columna
        \emph{Hash}). Abre el archivo \texttt{.sha256} con el Bloc de notas y
        compara: deben ser iguales. Basta con mirar los primeros y los
        últimos 6 caracteres \cite{ms-gethash}.
\end{pasos}

\subsection*{En Linux (otra distribución, o Lázaro ya instalado)}
Abre una terminal en la carpeta de la descarga y escribe:
\begin{terminal}
sha256sum -c lazaro-linux-*.sha256
\end{terminal}
Debe decir \texttt{La suma coincide} o \texttt{OK}.

\begin{cuidado}
Si los números \textbf{no coinciden}, la descarga está dañada: bórrala y
descárgala de nuevo, mejor desde la otra página (GitHub o SourceForge). No
sigas con una ISO dañada: el equipo podría no arrancar o fallar a mitad de
la instalación.
\end{cuidado}

\section{Prepara la memoria USB}\label{sec:usb}

Grabar la ISO en el USB no es lo mismo que copiarla: hace falta un programa
que la convierta en una memoria \emph{arrancable}. Elige uno de estos tres,
todos gratuitos.

\begin{cuidado}
Esto borra todo lo que tenga la memoria USB.
\end{cuidado}

\subsection{balenaEtcher (Windows, Mac y Linux): el más sencillo}
\begin{pasos}
  \item Descárgalo de \url{https://etcher.balena.io} e instálalo
        \cite{etcher}.
  \item Ábrelo y pulsa \menu{Flash from file}. Elige la ISO de Lázaro.
  \item Pulsa \menu{Select target} y elige tu memoria USB. Fíjate bien en
        el tamaño para no elegir otro disco.
  \item Pulsa \menu{Flash!} y espera. Al final Etcher comprueba la copia
        solo.
\end{pasos}

\subsection{Rufus (Windows)}
\begin{pasos}
  \item Descárgalo de \url{https://rufus.ie/es/}: no necesita instalarse
        \cite{rufus}.
  \item En \emph{Dispositivo}, elige tu memoria USB.
  \item Pulsa \menu{Seleccionar} y elige la ISO de Lázaro.
  \item Pulsa \menu{Empezar}. Si pregunta cómo escribir la imagen, elige
        \textbf{Escribir en modo Imagen DD}. Ese modo copia la ISO tal cual
        y es el que funciona con Lázaro.
\end{pasos}

\subsection{Ventoy: varias ISO en un mismo USB}
Ventoy prepara el USB una sola vez; después basta con copiar los archivos
ISO dentro, como a cualquier pendrive, y al arrancar te deja elegir cuál
usar \cite{ventoy}. Es cómodo si quieres llevar Lázaro junto a otros
sistemas.

\begin{cuidado}
Con Ventoy, después de copiar la ISO usa siempre \menu{Quitar el
dispositivo de forma segura} antes de sacar el USB. Si lo sacas de golpe,
la copia puede quedar dañada y el equipo no arranca o se detiene con un
\enquote{error de entrada/salida}. Copia también el archivo \texttt{.sha256} al
USB: así puedes repetir la comprobación sobre la copia que quedó en él.
\end{cuidado}

\section{BIOS, UEFI y el arranque seguro}\label{sec:bios-uefi}

Cuando se enciende, todo equipo ejecuta primero un pequeño programa propio
que revisa el hardware y busca desde dónde arrancar. En los equipos de
antes de 2012, aproximadamente, se llama \termino{BIOS}. En los más nuevos
es \termino{UEFI}, su reemplazo moderno \cite{uefi-org}. Lázaro funciona con
los dos y no tienes que hacer nada especial.

Muchos equipos con UEFI traen activado el \termino{arranque seguro}
(\emph{Secure Boot}), que solo deja arrancar sistemas firmados con llaves
de Microsoft \cite{ms-secureboot}. Si el equipo se niega a arrancar desde
el USB con un mensaje como \enquote{Secure Boot violation} o \enquote{imagen no
autorizada}, hay que desactivarlo:

\begin{pasos}
  \item Apaga el equipo. Enciéndelo y pulsa varias veces la tecla de la
        configuración (suele ser \keys{F2}, \keys{Supr} o \keys{Esc}; ver el
        apéndice~\ref{cap:b-teclas-arranque}).
  \item Busca la pestaña \emph{Security} o \emph{Boot} y la opción
        \emph{Secure Boot}. Ponla en \emph{Disabled}.
  \item Guarda y sal con \keys{F10} y confirma con \emph{Yes}.
\end{pasos}

\begin{nota}
Desactivar el arranque seguro no daña el equipo ni borra nada. Si más
adelante vuelves a Windows, puedes activarlo otra vez.
\end{nota}

\section{Arranca el equipo desde el USB}\label{sec:arrancar-usb}

\begin{pasos}
  \item Con el equipo apagado, conecta la memoria USB. En una PC de
        escritorio, mejor en un puerto de atrás.
  \item Enciéndelo y, apenas aparezca el logo de la marca, pulsa
        repetidamente la tecla del \textbf{menú de arranque}. Las más comunes
        están en la tabla~\ref{tab:teclas-cortas}; la lista completa, en el
        apéndice~\ref{cap:b-teclas-arranque}.
  \item En la lista que aparece, elige el USB con las flechas (puede
        aparecer con su marca, como \emph{SanDisk} o \emph{Kingston}, o como
        \emph{USB HDD}) y pulsa \keys{Enter}. Si aparece dos veces, una con
        \emph{UEFI} delante, elige esa.
\end{pasos}

\begin{table}[htbp]\centering
\caption{Tecla del menú de arranque en las marcas más comunes}\label{tab:teclas-cortas}
\begin{tabular}{ll}
\toprule
Marca & Tecla habitual \\
\midrule
Acer, Dell, Lenovo, Toshiba & \keys{F12} \\
Asus & \keys{Esc} o \keys{F8} \\
HP y Compaq & \keys{Esc} y luego \keys{F9} \\
Sony Vaio & \keys{F11} o el botón \emph{Assist} \\
PC de escritorio & \keys{F8}, \keys{F11} o \keys{F12} \\
\bottomrule
\end{tabular}
\end{table}

Si en la lista no aparece el USB, entra a la configuración del equipo
(sección anterior), busca la sección \emph{Boot} y pon el USB primero en
el orden de arranque.

\subsection{El menú de Lázaro}
Si todo va bien, aparece un menú con dos opciones (escritas sin tildes,
porque ese menú no las muestra bien):
\begin{description}
  \item[Lazaro Linux] la normal. Se elige sola a los pocos segundos.
  \item[Lazaro Linux (modo grafico seguro)] úsala solo si con la primera la
    pantalla se queda en negro. Arranca con una imagen más básica que
    funciona con casi cualquier tarjeta de video.
\end{description}

\captura{menu-arranque}{El menú de arranque de Lázaro en un equipo con UEFI}

Si el USB no arranca, o se detiene con un error, revisa la
sección~\ref{sec:prob-no-arranca} de \emph{Solución de problemas}.

\fuentes

\chapter{La sesión de prueba}\label{cap:03-sesion-de-prueba}

\begin{nota}
Capítulo en preparación.
\end{nota}

\chapter{Instalar paso a paso}\label{cap:04-instalar}

\begin{nota}
Capítulo en preparación.
\end{nota}

\chapter{El primer arranque}\label{cap:05-primer-arranque}

\section{Qué vas a ver}

Al encender aparece el logo de Lázaro con el pulso y, después, la pantalla
para entrar: elige tu usuario y escribe la contraseña. Si marcaste
\emph{Iniciar sesión automáticamente}, entra solo al escritorio.

\begin{nota}
Si en el disco también está Windows, al encender aparece primero un menú
negro, llamado \termino{GRUB}. La primera línea arranca Lázaro (se elige sola
a los pocos segundos); para Windows, baja con las flechas hasta
\emph{Windows Boot Manager} y pulsa \keys{Enter}.
\end{nota}

\section{Lo que pasa solo en los primeros minutos}

\begin{itemize}
  \item \textbf{La hora.} Si aparece corrida, se corrige sola en unos minutos
        al conectarse a internet.
  \item \textbf{Las aplicaciones de Flathub} que marcaste (OnlyOffice,
        Chrome, Brave) se descargan en segundo plano cuando hay internet.
        Dales unos minutos; aparecerán en el menú.
  \item \textbf{Las herramientas de Lázaro} se conectan a su repositorio para
        recibir mejoras con las actualizaciones normales
        (capítulo~\ref{cap:15-herramientas}).
\end{itemize}

\section{Si instalaste la Ligera sin internet}\label{sec:primer-arranque-texto}

Con la edición Ligera, una parte de la configuración final necesita
internet. Si instalaste sin conexión, el equipo arranca en una
\textbf{pantalla de texto} con un aviso. No es un error:

\begin{pasos}
  \item Conecta el cable de red, si tienes. Si no, entra con tu usuario y
        contraseña y conecta el wifi escribiendo:
\begin{terminal}
nmcli device wifi connect "NOMBRE DE LA RED" --ask
\end{terminal}
        Te pide la clave del wifi. Mientras la escribes no se ve nada en la
        pantalla: es normal.
  \item Espera. En unos minutos la configuración termina sola y aparece la
        pantalla de inicio. No hace falta reiniciar.
\end{pasos}

\section{Primeros pasos recomendados}

\begin{pasos}
  \item \textbf{Conecta el wifi} desde el icono de red de la barra
        (capítulo~\ref{cap:07-internet}). Se recuerda para siempre.
  \item \textbf{Actualiza.} Si aparece el icono de actualizaciones en la
        barra, haz clic y sigue los pasos (sección~\ref{sec:actualizar}).
  \item \textbf{Prueba la impresora}, si tienes (capítulo~\ref{cap:12-imprimir}).
  \item \textbf{Trae tus archivos} de la copia que hiciste
        (capítulo~\ref{cap:11-archivos}).
  \item \textbf{Abre Firefox} e inicia sesión en tu correo y tus cuentas.
  \item \textbf{Hojea la Parte II} de este manual para conocer lo que
        trae Lázaro.
\end{pasos}

El \textbf{Manual de Lázaro} queda siempre en el menú de aplicaciones, por
si lo necesitas.

\fuentes


\part{Uso diario}
\chapter{Tu escritorio}\label{cap:06-escritorio}

El \termino{escritorio} es lo que ves después de entrar: el fondo, la barra
con el menú y el reloj, y las ventanas de los programas. Este capítulo
explica XFCE, el que viene marcado de fábrica, y al final resume las
diferencias de los demás.

\section{XFCE de un vistazo}\label{sec:xfce}

\captura{xfce-escritorio}{El escritorio XFCE de Lázaro}

XFCE trae dos barras, que XFCE llama \termino{paneles} \cite{xfce-docs}:

\begin{description}
  \item[El panel de arriba.] A la izquierda, el menú \menu{Aplicaciones}
    con todos los programas, ordenados por tipo (Internet, Oficina,
    Multimedia, Sistema\ldots). En el medio, un botón por cada ventana
    abierta: haz clic para ir a ella o para minimizarla. A la derecha, la
    \termino{bandeja}: los iconos de la red, el volumen, la batería, el
    aviso de actualizaciones y el reloj.
  \item[El panel de abajo.] Un pequeño \enquote{muelle} con accesos a los
    programas más usados: mostrar el escritorio, la terminal, la carpeta
    personal, el navegador y la búsqueda de programas.
\end{description}

\subsection{Abrir un programa}
\begin{pasos}
  \item Haz clic en \menu{Aplicaciones}, arriba a la izquierda.
  \item Pasa el ratón por una categoría, por ejemplo \menu{Oficina}.
  \item Haz clic en el programa.
\end{pasos}

\begin{consejo}
Más rápido: pulsa \keys{Alt + F2} (o la lupa del panel de abajo), escribe
las primeras letras del programa, por ejemplo \texttt{fire} para Firefox, y
pulsa \keys{Enter}.
\end{consejo}

\subsection{Las ventanas}
Cada ventana tiene arriba a la derecha tres botones: \emph{minimizar}
(la esconde en el panel), \emph{maximizar} (ocupa toda la pantalla) y
\emph{cerrar}. Para moverla, arrástrala desde la barra del título. Para
cambiar su tamaño, arrastra un borde o una esquina. Para pasar de una
ventana a otra con el teclado: \keys{Alt + Tab}.

\subsection{Poner un programa en el panel o en el escritorio}
Abre el menú \menu{Aplicaciones}, haz clic derecho sobre el programa y
elige \menu{Añadir al panel} o \menu{Añadir al escritorio}.

\subsection{Cambiar el fondo, la letra y los colores}
\begin{itemize}
  \item \textbf{Fondo:} clic derecho en el escritorio, \menu{Configuración del
        escritorio}. Elige una imagen; para usar una foto tuya, cambia la
        carpeta de arriba.
  \item \textbf{Tamaño de letra y tema claro u oscuro:}
        \menu{Aplicaciones > Configuración > Apariencia}. En la pestaña
        \menu{Estilo} eliges el tema; en \menu{Tipos de letra}, el tamaño.
  \item \textbf{Todo lo demás:} \menu{Aplicaciones > Configuración >
        Gestor de configuración} reúne todos los ajustes en una ventana
        (capítulo~\ref{cap:14-ajustes}).
\end{itemize}

\begin{consejo}
¿Te cuesta leer? Sube el tamaño de letra a 12 o 13 en \menu{Apariencia >
Tipos de letra}, y en Firefox usa \keys{\ctrl + +} para agrandar la página.
\end{consejo}

\subsection{Apagar, reiniciar y cerrar sesión}
En el menú \menu{Aplicaciones}, abajo, está \menu{Salir}. Ofrece
\emph{Apagar}, \emph{Reiniciar}, \emph{Cerrar la sesión} y, en portátiles,
\emph{Suspender}. Apaga siempre así: cortar la corriente de golpe puede
dañar archivos que se estaban guardando.

\section{Los otros escritorios}

Si elegiste otro escritorio al instalar, todo lo de este libro sigue
valiendo: los programas son los mismos. Cambian el aspecto y los nombres
de algunos menús.

\begin{description}
  \item[LXQt] Muy parecido a XFCE pero con una sola barra, abajo, como en
    Windows. El menú está en el botón de la izquierda; la configuración, en
    \menu{Preferencias > Configuración de LXQt} \cite{lxqt-wiki}. Su
    gestor de archivos se llama \emph{PCManFM-Qt}.
  \item[MATE] Dos paneles como XFCE. Arriba, los menús \menu{Aplicaciones},
    \menu{Lugares} y \menu{Sistema}; la configuración está en
    \menu{Sistema > Preferencias} \cite{mate-docs}. Su gestor de archivos se
    llama \emph{Caja}.
  \item[KDE Plasma] Una barra abajo, como Windows, y un menú que se abre con
    el botón de la izquierda o con la tecla \keys{Windows}. Todo se ajusta
    en \emph{Preferencias del sistema}. Los archivos se manejan con
    \emph{Dolphin} \cite{kde-userbase}.
  \item[GNOME] Una barra fina arriba. Pulsa \keys{Windows} (o
    \emph{Actividades}, arriba a la izquierda) para ver las ventanas
    abiertas y buscar programas escribiendo su nombre. La configuración
    está en \emph{Configuración}; los archivos, en \emph{Archivos}
    \cite{gnome-help}.
  \item[Sway] No tiene menús: todo se hace con el teclado. Las ventanas
    se acomodan solas una al lado de la otra. Lo básico: \keys{Windows +
    Enter} abre una terminal, \keys{Windows + D} busca programas,
    \keys{Windows + Mayús + Q} cierra la ventana. Recomendado solo para
    quien ya lo conoce \cite{sway-wiki}.
  \item[XFCE Wayland] Se ve y se usa igual que XFCE. Usa una tecnología
    gráfica más nueva, todavía en prueba: si algo falla, XFCE es la opción
    segura.
\end{description}

\begin{nota}
En los escritorios KDE Plasma y LXQt la tienda de aplicaciones se llama
\emph{Discover}; en los demás, \emph{Software} (capítulo~\ref{cap:13-programas}).
\end{nota}

\fuentes

\chapter{Internet}\label{cap:07-internet}

\begin{nota}
Capítulo en preparación.
\end{nota}

\chapter{Documentos}\label{cap:08-documentos}

\begin{nota}
Capítulo en preparación.
\end{nota}

\chapter{Archivos PDF}\label{cap:09-pdf}

\begin{nota}
Capítulo en preparación.
\end{nota}

\chapter{Videos, música y fotos}\label{cap:10-multimedia}

\section{YouTube ligero}\label{sec:youtube-ligero}

Ver YouTube en el navegador gasta mucho procesador: en equipos viejos el
video se traba aunque la conexión sea buena. \textbf{YouTube ligero}
busca los videos en una ventana sencilla (\emph{Pipe Viewer}) y los
reproduce con \emph{mpv}, un reproductor muy liviano, en 480p, el formato
que menos procesador gasta \cite{mpv,ytdlp}. No necesita navegador ni
cuenta.

\begin{pasos}
  \item Abre \menu{Aplicaciones > Multimedia > YouTube}.
  \item Escribe lo que buscas arriba y pulsa \keys{Enter}.
  \item Haz doble clic en un video. Se abre en una ventana de mpv.
\end{pasos}

\captura{youtube-ligero}{YouTube ligero: la búsqueda en Pipe Viewer}

Mientras se reproduce: \keys{Espacio} pausa, \keys{\arrowkeyleft} y
\keys{\arrowkeyright} retroceden y adelantan 5 segundos, \keys{F} pone
pantalla completa, \keys{9} y \keys{0} bajan y suben el volumen, \keys{Q}
cierra.

\begin{description}
  \item[Si el video se traba] (procesadores de un solo núcleo): en Pipe
    Viewer, abre el menú de preferencias y elige resolución \emph{360}. Se
    guarda al cerrar la ventana.
  \item[Si deja de funcionar] YouTube cambia seguido y a veces rompe los
    programas que lo leen. Instala las actualizaciones
    (sección~\ref{sec:actualizar}); si no alcanza, abre una terminal y
    escribe:
\begin{terminal}
doas apk upgrade yt-dlp yt-dlp-ejs pipe-viewer
\end{terminal}
\end{description}

\section{Videos y música}

Lázaro abre los videos y los audios con \textbf{mpv}: haz doble clic en el
archivo. Reproduce casi cualquier formato, incluidos los audios de
WhatsApp. Las teclas son las mismas de YouTube ligero.

Si al instalar marcaste \textbf{VLC}, tienes un reproductor más completo,
con menús, listas de reproducción y subtítulos \cite{vlc}. Para abrir un
archivo con VLC: clic derecho, \menu{Abrir con > VLC}. Para que sea el
predeterminado, en esa misma ventana marca la casilla para usarlo siempre.

\begin{consejo}
¿Una película con subtítulos en un archivo aparte (\texttt{.srt})? Ponle el
mismo nombre que al video y déjalo en la misma carpeta: mpv y VLC lo
cargan solos.
\end{consejo}

\subsection{El sonido}
El icono del altavoz en la bandeja sube y baja el volumen; con clic
derecho (o \menu{Configuración del sonido}) se elige la salida: los
parlantes, los audífonos, la pantalla por HDMI o unos audífonos
Bluetooth. Si no se escucha nada, mira la sección~\ref{sec:prob-sonido}.

\section{Fotos}

Haz doble clic en una foto para verla. Con las flechas pasas a la
siguiente. En XFCE se abre con \emph{Ristretto}; en los demás escritorios,
con su visor propio.

\begin{description}
  \item[Pasar las fotos del celular] Conecta el celular por cable. En
    Android, desliza la barra de arriba, toca la notificación del USB y
    elige \emph{Transferencia de archivos}. El celular aparece en el
    gestor de archivos como un disco más (capítulo~\ref{cap:11-archivos}).
  \item[Retocar] Para recortar, girar o mejorar el color, la tienda
    ofrece programas como \emph{GIMP} (completo) o \emph{Pinta} (sencillo).
\end{description}

\section{Capturas de pantalla}

En XFCE, la tecla \keys{ImprPant} abre \emph{Captura de pantalla}: elige si
quieres toda la pantalla, una ventana o un área, y dónde guardarla. En
GNOME y KDE Plasma esa misma tecla abre su herramienta propia.

\section{Cámara web}

Para probar la cámara, abre una videollamada de prueba en el navegador
(por ejemplo, \url{https://meet.jit.si}). Si la página pide permiso para
usar la cámara, pulsa \menu{Permitir}.

\fuentes

\chapter{Archivos, pendrives y copias de seguridad}\label{cap:11-archivos}

\section{El gestor de archivos}

Tus archivos se manejan con el \termino{gestor de archivos}: en XFCE se
llama \emph{Thunar}; en otros escritorios, Caja (MATE), PCManFM-Qt
(LXQt), Archivos (GNOME) o Dolphin (KDE Plasma). Se abre desde el panel de
abajo, con el icono de la carpeta, o haciendo doble clic en \emph{Carpeta
personal} del escritorio.

\captura{thunar}{Thunar, el gestor de archivos de XFCE}

\subsection{Tu carpeta personal}
Cada usuario tiene su \emph{carpeta personal}, el equivalente a la carpeta
de usuario de Windows. Dentro están \emph{Documentos}, \emph{Descargas},
\emph{Imágenes}, \emph{Música}, \emph{Videos} y \emph{Escritorio}. Guarda
ahí lo tuyo: el resto del disco es del sistema y no hace falta tocarlo.

\subsection{Lo de todos los días}
\begin{itemize}
  \item \textbf{Copiar}: clic derecho sobre el archivo, \menu{Copiar}; ve a
        la otra carpeta, clic derecho en un espacio vacío, \menu{Pegar}.
        También se puede arrastrar.
  \item \textbf{Cambiar el nombre}: clic sobre el archivo y \keys{F2}.
  \item \textbf{Borrar}: \keys{Supr} lo manda a la \emph{Papelera}; desde
        ahí se puede recuperar hasta que la vacías.
  \item \textbf{Crear una carpeta}: clic derecho en un espacio vacío,
        \menu{Crear carpeta}.
  \item \textbf{Buscar}: \keys{\ctrl + F} y escribe parte del nombre.
  \item \textbf{Ver los archivos ocultos}: \keys{\ctrl + H}. Los archivos
        cuyo nombre empieza con un punto son configuraciones; mejor no
        tocarlos.
\end{itemize}

\subsection{Archivos comprimidos (.zip, .rar, .7z)}
Clic derecho sobre el archivo, \menu{Extraer aquí}. Para comprimir:
selecciona los archivos, clic derecho, \menu{Crear un archivador} (o
\menu{Comprimir}) y elige el formato \texttt{.zip}, el que se abre en
cualquier equipo.

\section{Pendrives, discos externos y el celular}

Al conectar un pendrive o un disco externo aparece solo en la columna de
la izquierda del gestor de archivos. Haz clic para abrirlo.

\begin{cuidado}
Antes de sacar el pendrive, haz clic en el triángulo de expulsar que
aparece al lado de su nombre (o clic derecho, \menu{Expulsar}) y espera a
que desaparezca. Si lo sacas mientras se está copiando algo, el archivo
y a veces todo el pendrive pueden quedar dañados.
\end{cuidado}

\begin{nota}
Copiar archivos grandes a un pendrive lento puede tardar: el sistema
copia primero a la memoria y después al pendrive. El Afinador de Lázaro
ajusta esto para que el equipo no se congele mientras tanto
(capítulo~\ref{cap:15-herramientas}). Espera siempre a expulsar.
\end{nota}

\subsection{El celular}
Conecta el celular por cable y, en Android, elige \emph{Transferencia de
archivos} en la notificación del USB. Aparece en el gestor de archivos y
puedes copiar las fotos de la carpeta \texttt{DCIM}. Con un iPhone,
desbloquéalo y toca \emph{Confiar} cuando lo pregunte.

\subsection{Formatear un pendrive}
Formatear borra todo y lo deja como nuevo. En el gestor de archivos, clic
derecho sobre el pendrive, \menu{Formatear} (si tu escritorio no lo
ofrece, instala \emph{Discos} desde la tienda). Elige \textbf{exFAT} si
lo vas a usar también en Windows, Mac o un televisor con archivos de más
de 4 GB; \textbf{FAT} para lo más compatible con equipos antiguos.

\section{Copias de seguridad}\label{sec:copias}

Un disco puede fallar sin avisar, sobre todo en un equipo viejo. La regla
que recomiendan los especialistas es simple: tener \textbf{al menos dos
copias} de lo importante, en lugares distintos \cite{incibe-copias}.

\begin{pasos}
  \item Consigue un disco externo o un pendrive grande que uses solo para
        esto.
  \item Una vez por semana (o cada vez que guardes algo importante),
        conéctalo y copia tus carpetas \emph{Documentos} e \emph{Imágenes}
        dentro de una carpeta con la fecha, por ejemplo
        \texttt{Copia 2026-11-09}.
  \item Expúlsalo y guárdalo en otro lugar.
\end{pasos}

\begin{consejo}
Si prefieres que se haga sola, la tienda de aplicaciones ofrece
\emph{Déjà Dup} (\enquote{Copias de seguridad}), que guarda una copia
programada en un disco externo o en la nube y permite recuperar versiones
anteriores de un archivo.
\end{consejo}

\fuentes

\chapter{Imprimir y escanear}\label{cap:12-imprimir}

Lázaro usa \termino{CUPS}, el sistema de impresión de Linux y de macOS
\cite{cups}. Al instalar revisa las impresoras conectadas y deja puestos
los controladores de las marcas más comunes: HP, Epson, Brother, Samsung y
Xerox.

\section{Agregar la impresora}

\subsection{Por cable USB}
Conéctala y enciéndela. En la mayoría de los casos se agrega sola en
unos segundos. Para comprobarlo, imprime cualquier cosa: la impresora
debe aparecer en la lista.

\subsection{Por wifi o red}
Las impresoras modernas que están en la misma red wifi se detectan solas
y aparecen en la lista al imprimir. Si no aparece:
\begin{pasos}
  \item Abre \menu{Aplicaciones > Sistema > Configuración de impresión}.
  \item Pulsa \menu{Añadir} y escribe tu contraseña.
  \item Elige la impresora en la lista \emph{Impresora de red}. Si no
        está, escribe la dirección IP que muestra la propia impresora en su
        pantalla o en la hoja de configuración que imprime.
  \item Elige la marca y el modelo, y pulsa \menu{Aplicar}.
  \item Imprime una página de prueba.
\end{pasos}

\captura{impresoras}{La ventana de Configuración de impresión}

\begin{consejo}
Si conectas una impresora después de instalar y le falta su controlador,
aparece el aviso \enquote{Falta algo para tu hardware}: pulsa
\menu{Instalar} (sección~\ref{sec:controladores}). ¿Tu modelo funciona en
Linux? La base de datos de OpenPrinting lo dice \cite{openprinting}.
\end{consejo}

\section{Imprimir}

En cualquier programa: \menu{Archivo > Imprimir} o \keys{\ctrl + P}.
Revisa la impresora elegida, la cantidad de copias y las páginas (por
ejemplo, \texttt{1-3,5}). En \emph{Propiedades} o \emph{Configuración de
página} se elige doble cara, blanco y negro o la calidad.

\section{Escanear}

Lázaro instala el programa \emph{Escáner de documentos} (Simple Scan) y
los controladores de escáner de Linux, llamados \termino{SANE}
\cite{sane}. Funciona con la mayoría de los escáneres USB y de las
impresoras multifunción en red.

\begin{pasos}
  \item Abre \menu{Aplicaciones > Gráficos > Escáner de documentos}.
  \item Pon la hoja y pulsa \menu{Escanear}. Elige arriba \emph{Texto} para
        documentos o \emph{Foto} para fotografías.
  \item Para varias hojas, pulsa \menu{Escanear} otra vez: se agregan como
        páginas nuevas.
  \item Pulsa \menu{Guardar}: elige PDF para documentos, JPG para fotos.
\end{pasos}

Las multifunción HP necesitan además su paquete \emph{hplip}, que Lázaro
instala cuando detecta una impresora HP.

\fuentes

\chapter{Instalar, quitar y actualizar programas}\label{cap:13-programas}

\begin{nota}
Capítulo en preparación.
\end{nota}

\chapter{Ajustes del equipo}\label{cap:14-ajustes}

En XFCE todos los ajustes están en \menu{Aplicaciones > Configuración >
Gestor de configuración}. En los demás escritorios: \menu{Sistema >
Preferencias} (MATE), \menu{Preferencias > Configuración de LXQt} (LXQt),
\emph{Configuración} (GNOME) o \emph{Preferencias del sistema} (KDE
Plasma).

\captura{gestor-configuracion}{El Gestor de configuración de XFCE}

\section{Pantalla}

\menu{Configuración > Pantalla} cambia la resolución (el tamaño de todo en
la pantalla) y, si conectas un monitor o un proyector, si se ve lo mismo
en los dos o si se extiende el escritorio.

\begin{consejo}
Si todo se ve demasiado pequeño en una pantalla grande, en lugar de bajar
la resolución sube el tamaño de letra en \menu{Apariencia > Tipos de
letra} y, en \menu{Apariencia > Configuración}, la \emph{escala}.
\end{consejo}

\section{Sonido}

El icono del altavoz de la bandeja sube y baja el volumen. Con clic
derecho, \menu{Configuración del sonido} muestra todas las salidas y
entradas: elige ahí los audífonos, el parlante Bluetooth o el micrófono
para las videollamadas.

\section{Teclado y ratón}

\begin{itemize}
  \item \menu{Configuración > Teclado > Distribución} cambia el idioma del
        teclado. Para usar dos (por ejemplo, español e inglés), agrégalos
        aquí y elige una combinación para cambiar entre ellos.
  \item \menu{Configuración > Ratón y panel táctil}: velocidad del puntero,
        \emph{tocar para hacer clic} y el desplazamiento con dos dedos en el
        touchpad.
\end{itemize}

\section{Fecha y hora}

La hora se ajusta sola por internet. Para cambiar la zona horaria, haz
clic derecho en el reloj del panel, \menu{Propiedades}. Si la hora está mal
después de usar Windows en el mismo equipo, mira la
sección~\ref{sec:prob-hora}.

\section{Idioma}

Lázaro queda en el idioma que elegiste al instalar, y todos los programas
lo siguen. Cambiar el idioma de un sistema ya instalado todavía no tiene un
botón: si lo necesitas, pide ayuda (capítulo~\ref{cap:17-pedir-ayuda}).
Para escribir en otro idioma no hace falta cambiarlo: basta con agregar
otra distribución de teclado (arriba) y, en LibreOffice, elegir el idioma
del texto en \menu{Herramientas > Idioma}.

\section{Usuarios y contraseñas}

\subsection{Cambiar tu contraseña}
Abre una terminal, escribe \texttt{passwd} y pulsa \keys{Enter}. Pide la
contraseña actual y dos veces la nueva. Mientras escribes no se ve nada: es
normal.

\subsection{Entrar sin contraseña al encender}
Lázaro recuerda esta opción en el archivo \texttt{/etc/lazaro/autologin}.
Para dejar de entrar solo, en la terminal:
\begin{terminal}
doas rm /etc/lazaro/autologin
\end{terminal}
Para empezar a entrar solo con tu usuario:
\begin{terminal}
echo "\$USER" | doas tee /etc/lazaro/autologin
\end{terminal}
En los dos casos, reinicia para que tome efecto.

\subsection{Otra persona en el mismo equipo}
Cada persona puede tener su propio usuario, con su carpeta y su
configuración. En la terminal, cambiando \texttt{ana} por el nombre:
\begin{terminal}
doas adduser ana
\end{terminal}

\subsection{Qué es doas}
Algunas órdenes cambian el sistema y necesitan permiso de administrador.
Se escriben con \texttt{doas} delante, que pide tu contraseña. Es lo mismo
que el \emph{sudo} de otras distribuciones o el \enquote{Ejecutar como
administrador} de Windows. Usa \texttt{doas} solo con órdenes que
entiendas o que vengan de una fuente de confianza, como este manual.

\section{Energía y batería}

\menu{Configuración > Gestor de energía} decide cuándo se apaga la
pantalla, qué pasa al cerrar la tapa del portátil y a qué nivel de batería
se suspende. Si el equipo tiene batería, Lázaro instala los \emph{modos de
energía}: en el icono de la batería se elige entre \emph{ahorro},
\emph{equilibrado} y \emph{rendimiento}.

\section{Bluetooth}

En KDE Plasma y GNOME, el Bluetooth está en la configuración del sistema y
en el menú de arriba a la derecha (GNOME) o en la bandeja (Plasma). En
XFCE, MATE y LXQt, instala desde la tienda el \emph{Gestor de Bluetooth}
(Blueman): agrega un icono a la bandeja para buscar y emparejar
audífonos, parlantes, ratones y teclados.

Para emparejar, pon el aparato en modo de emparejamiento (casi siempre,
manteniendo apretado su botón de encendido hasta que parpadee) y elígelo
en la lista.

\fuentes


\part{Herramientas de Lázaro}
\chapter{Las herramientas de Lázaro}\label{cap:15-herramientas}

\begin{nota}
Capítulo en preparación.
\end{nota}


\part{Cuando algo falla}
\chapter{Solución de problemas}\label{cap:16-problemas}

\begin{nota}
Capítulo en preparación.
\end{nota}

\chapter{Pedir ayuda}\label{cap:17-pedir-ayuda}

\begin{nota}
Capítulo en preparación.
\end{nota}

\chapter{La terminal sin miedo}\label{cap:18-terminal}

La \termino{terminal} es una ventana donde se le dan órdenes al equipo
escribiendo, en lugar de hacer clic. Para casi todo lo de este libro no
hace falta, pero a veces es el camino más corto, y es como suelen
responder quienes te ayudan.

Se abre desde \menu{Aplicaciones > Sistema > Terminal} (o el icono del
panel de abajo en XFCE).

\captura{terminal}{La terminal de XFCE}

\section{Cómo se usa}

\begin{itemize}
  \item Escribe la orden y pulsa \keys{Enter}.
  \item Las mayúsculas, las minúsculas y los espacios importan.
  \item \keys{\arrowkeyup} trae la orden anterior; así no hay que escribirla
        de nuevo.
  \item \keys{Tab} completa nombres de archivos y órdenes: escribe las
        primeras letras y pulsa \keys{Tab}.
  \item Para pegar, usa \keys{\ctrl + Mayús + V} o clic derecho,
        \menu{Pegar} (\keys{\ctrl + V} solo no funciona aquí).
  \item \keys{\ctrl + C} detiene la orden que está corriendo.
  \item Cuando una orden pide la contraseña, lo que escribes no se ve.
        Escríbela igual y pulsa \keys{Enter}.
\end{itemize}

\section{Órdenes útiles}

\begin{center}
\begin{tabularx}{\textwidth}{>{\ttfamily}lX}
\toprule
\normalfont Orden & Qué hace \\
\midrule
ls & Muestra los archivos de la carpeta actual \\
cd Documentos & Entra en la carpeta \emph{Documentos} \\
cd & Vuelve a la carpeta personal \\
pwd & Dice en qué carpeta estás \\
fastfetch & Resumen del equipo y del sistema \\
btop & Qué está usando el procesador y la memoria (\keys{Q} para salir) \\
df -h & Cuánto espacio queda en los discos \\
free -h & Cuánta memoria está en uso \\
ip addr & Las conexiones de red y sus direcciones \\
nmcli device wifi list & Las redes wifi que se ven \\
doas apk upgrade & Actualiza todo lo de Alpine \\
lazaro-diagnostico & El informe para pedir ayuda \\
\bottomrule
\end{tabularx}
\end{center}

\section{Ser administrador: doas}

Las órdenes que cambian el sistema se escriben con \texttt{doas} delante:
\begin{terminal}
doas apk add gimp
\end{terminal}
\texttt{doas} pide tu contraseña y ejecuta la orden con permisos de
administrador. Es el equivalente a \enquote{Ejecutar como administrador} de
Windows o al \texttt{sudo} de otras distribuciones. Úsalo solo con órdenes
que entiendas o que vengan de una fuente de confianza.

\begin{nota}
Para aprender más, la wiki de Alpine explica cada parte del sistema
\cite{alpine-wiki}, y el gestor de paquetes \texttt{apk} tiene su propia
página \cite{alpine-apk}.
\end{nota}

\fuentes


\appendix
\chapter{Glosario}\label{cap:a-glosario}

\begin{nota}
Capítulo en preparación.
\end{nota}

\chapter{Teclas para arrancar desde el USB}\label{cap:b-teclas-arranque}

\begin{nota}
Capítulo en preparación.
\end{nota}

\chapter{Si venías de Windows}\label{cap:c-equivalencias}

\begin{center}
\begin{tabularx}{\textwidth}{XX}
\toprule
En Windows & En Lázaro \\
\midrule
Menú Inicio & Menú \menu{Aplicaciones} (arriba a la izquierda en XFCE) \\
Explorador de archivos & Gestor de archivos (Thunar en XFCE) \\
Panel de control / Configuración & \menu{Aplicaciones > Configuración >
  Gestor de configuración} \\
Microsoft Store & La tienda de aplicaciones (Software o Discover) \\
Windows Update & El aviso de actualizaciones, \menu{Actualizar el sistema} \\
Administrador de tareas & Btop (\texttt{btop} en la terminal) o el monitor
  del sistema de tu escritorio \\
Word, Excel, PowerPoint & LibreOffice Writer, Calc, Impress (u OnlyOffice) \\
Edge, Chrome & Firefox (Chrome y Brave, opcionales) \\
Adobe Reader & Visor de documentos (Evince) \\
Reproductor de Windows Media & mpv o VLC \\
Fotos & Ristretto (visor de fotos de XFCE) \\
Bloc de notas & Mousepad (editor de texto de XFCE) \\
Recortes / Imprimir pantalla & Captura de pantalla (\keys{ImprPant}) \\
Calculadora & Galculator \\
WinRAR, 7-Zip & Clic derecho, \menu{Extraer aquí} \\
\enquote{Ejecutar como administrador} & \texttt{doas} delante de la orden \\
Símbolo del sistema, PowerShell & La terminal \\
Antivirus & No hace falta: los programas vienen de fuentes firmadas y se
  actualizan solos. Mantén el sistema al día. \\
Unidad C: & La carpeta personal, para tus archivos \\
\bottomrule
\end{tabularx}
\end{center}

\chapter{Atajos de teclado}\label{cap:d-atajos}

Funcionan en casi todos los programas y escritorios.

\begin{center}
\begin{tabularx}{\textwidth}{lX}
\toprule
Teclas & Qué hacen \\
\midrule
\keys{\ctrl + C} / \keys{\ctrl + X} / \keys{\ctrl + V} & Copiar / cortar / pegar \\
\keys{\ctrl + Z} / \keys{\ctrl + Y} & Deshacer / rehacer \\
\keys{\ctrl + A} & Seleccionar todo \\
\keys{\ctrl + S} & Guardar \\
\keys{\ctrl + P} & Imprimir \\
\keys{\ctrl + F} & Buscar \\
\keys{\ctrl + +} / \keys{\ctrl + -} & Agrandar / achicar \\
\keys{Alt + Tab} & Pasar de una ventana a otra \\
\keys{Alt + F4} & Cerrar la ventana \\
\keys{F2} & Cambiar el nombre de un archivo \\
\keys{Supr} & Mandar a la papelera \\
\keys{ImprPant} & Captura de pantalla \\
\bottomrule
\end{tabularx}
\end{center}

\section*{En XFCE}
\begin{center}
\begin{tabularx}{\textwidth}{lX}
\toprule
\keys{Alt + F2} & Buscar y abrir un programa \\
\keys{\ctrl + Alt + T} & Abrir la terminal \\
\keys{\ctrl + Alt + D} & Mostrar el escritorio \\
\keys{\ctrl + Alt + Supr} & Bloquear la pantalla \\
\bottomrule
\end{tabularx}
\end{center}

\section*{En Firefox}
\begin{center}
\begin{tabularx}{\textwidth}{lX}
\toprule
\keys{\ctrl + T} / \keys{\ctrl + W} & Abrir / cerrar pestaña \\
\keys{\ctrl + Mayús + T} & Recuperar la pestaña cerrada \\
\keys{\ctrl + L} & Ir a la barra de direcciones \\
\keys{\ctrl + D} & Guardar en marcadores \\
\keys{F5} & Recargar la página \\
\keys{F11} & Pantalla completa \\
\bottomrule
\end{tabularx}
\end{center}

\backmatter
\printbibliography[heading=bibintoc,title={Bibliografía}]
\end{document}
